\input{../tex-inputs/latex-leseansicht-vorspann}

\section[Gustav Schwarzkopf an Arthur Schnitzler, {[}8.? 9. 1893{]}]{L04240 Gustav Schwarzkopf an Arthur Schnitzler, {[}8.? 9. 1893{]}}
\nopagebreak\mylabel{L04240v}
\rehead{ }\normalsize\beginnumbering\briefempfaengerindex{Schnitzler, Arthur@\textsc{Schnitzler, Arthur}!zzzSchwarzkopf, Gustav@\emph{von Gustav Schwarzkopf}!1893-09-081@{{[}8.? 9. 1893{]}}|(be}
\toendnotes[C]{\smallbreak\pagebreak[2]}
\correspDesc{Versand  durch Gustav Schwarzkopf am [8.? 9. 1893] in Wien
\newline{}Erhalt  durch Arthur Schnitzler am [8.? 9. 1893] in Wien}\toendnotes[C]{\smallbreak}
\Standort{CUL, Schnitzler, B 96a.}
\physDesc{Visitenkarte, 209 Zeichen
\newline{}Handschrift: schwarze Tinte, deutsche Kurrent}\toendnotes[C]{\smallbreak}
\pstart{}{\pb}Lieber Freund!\pend\vspace{0.5em}
\pstart
           Ich fahre \label{K_L04240-1v}\edtext{heute
                  Nachmittag in die Brühl\oindex{Brühl@\textbf{Brühl}, \emph{Tal}|pw} und
               hoffe Sie Abends zu{ }ſehen}{\lemma{\textnormal{\emph{heute … sehen}}}\Cendnote{\textnormal{Schnitzler war im »Sept 93« nur am 8. 9. 1893 in der Brühl\oindex{Brühl@\textbf{Brühl}, \emph{Tal}|pwk}. Er
                  traf dort auf Schwarzkopf\pwindex{Schwarzkopf, Gustav 7.\,11.\,1853 Wien – 13.\,11.\,1939 ebd.@\textsc{Schwarzkopf, Gustav} (7.\,11.\,1853 Wien – 13.\,11.\,1939 ebd.), \emph{Schriftsteller}|pwk}, womit eine
                  halbwegs verlässliche Datierung dieser Visitenkarte möglich wird.}}}\label{K_L04240-1}.
               Vielleicht überlegen Sie{ }ſich’s doch und \label{K_L04240-2v}\edtext{übernachten bei Schönberger\oindex{Zum goldenen Stern@\textbf{Zum goldenen Stern}, \emph{Lokal}|pw}}{\lemma{\textnormal{\emph{übernachten bei Schönberger}}}\Cendnote{\textnormal{Im \emph{Tagebuch}\pwindex{Schnitzler, Arthur 15.\,5.\,1862 Wien – 21.\,10.\,1931 ebd.@\textsc{Schnitzler, Arthur} (15.\,5.\,1862 Wien – 21.\,10.\,1931 ebd.), \emph{Schriftsteller, Mediziner}!Tagebuch@\strich\emph{Tagebuch}|pwk}-Eintrag zum 8. 9. 1893 erwähnt
                  Schnitzler, noch am selben Abend wieder nach Wien\oindex{Wien@\textbf{Wien}, \emph{Verwaltungsgebiet}|pwk} gefahren zu sein.}}}\label{K_L04240-2}, wir kö{\geminationn}ten dann eine
               hübſche Partie {\pb}machen.\pend
           
\pstart
           Herzlich grüßend{\\[\baselineskip]} Ihr{\\[\baselineskip]}\pend
           \leftskip=0em{}
\pstart
           \centering{}\textcolor{gray}{\textbf{Gustav Schwarzkopf}}\pend
           \selectlanguage{ngerman}\endnumbering\briefempfaengerindex{Schnitzler, Arthur@\textsc{Schnitzler, Arthur}!zzzSchwarzkopf, Gustav@\emph{von Gustav Schwarzkopf}!1893-09-081@{{[}8.? 9. 1893{]}}|)be}\mylabel{L04240h}
\begin{anhang}
\end{anhang}\newcommand{\dateiname}{L04240}\newcommand{\titel}{Gustav Schwarzkopf an Arthur Schnitzler, [8.? 9. 1893]}\newcommand{\editorInnen}{Herausgegeben von Jahnke, SelmaMüller, Martin Anton}\input{../tex-inputs/latex-leseansicht-abspann}
